\documentclass[unicode]{beamer}
\linespread{1.3}
\usepackage[]{cite}
\usepackage{cmap}
\usepackage[T2A]{fontenc}
\usepackage[utf8]{inputenc}
\usepackage[english, russian]{babel}
\usepackage{amsmath,mathrsfs,amsfonts,amssymb}
\usepackage{graphicx, epsfig}
\usepackage{subfig}
\usepackage{color}
\usepackage{wrapfig}
\usepackage{setspace}
\usepackage{dsfont}
\usepackage{afterpage} 
\usepackage{mathtools}
\usepackage{bm}
\usepackage{NotationStyle}
\usepackage{tabularx}
\usepackage{array}
\usepackage{makecell}
\usepackage{multirow}
\usepackage{rotating}
\usepackage{rotate}
\usepackage{comment}
\usepackage{float}
\usepackage{upgreek}
\usepackage{graphicx}
\usepackage{amsmath}
\usepackage{amsthm}
\usepackage{amssymb}
\usepackage{booktabs}
\usepackage{tabularx}
\usepackage{multirow}
\usepackage{hyperref}
\usepackage{algorithm}
\usepackage{algorithmic}
\usepackage{xcolor}
\usepackage{tikz}
\usetikzlibrary{arrows,positioning,shapes.geometric}
\parindent=1.25cm
%\usepackage[left=3cm, right=1.5cm, top=2cm, bottom=2cm, bindingoffset=0cm]{geometry}

%\newcommand{\vect}[1]{\boldsymbol{\mathbf{#1}}}
\newcommand{\vect}[1]{\textbf{#1}}

\usepackage[shortlabels]{enumitem}
                    \setlist[enumerate, 1]{1\textsuperscript{o}}

\usepackage{caption}

\DeclareCaptionStyle{mystyle}
{format=plain,%
	textformat=period,
	justification=RaggedRight,
	singlelinecheck=true,
}

\DeclareCaptionStyle{singlelineraggedleft}
[justification=RaggedLeft]% right aligned if single line and no `singlelinecheck=false`
{style=mystyle}

\newcommand\argmin{\mathop{\arg\min}}
\newcommand{\T}{^{\text{\tiny\sffamily\upshape\mdseries T}}}
\newcommand{\hchi}{\hat{\boldsymbol{\chi}}}
\newcommand{\hphi}{\hat{\boldsymbol{\varphi}}}
\newcommand{\bchi}{\boldsymbol{\chi}}
\newcommand{\hx}{\hat{x}}
\newcommand{\hy}{\hat{y}}
\newcommand{\M}{\mathcal{M}}
\newcommand{\N}{\mathcal{N}}
\newcommand{\p}{p(\cdot)}
\newcommand{\q}{q(\cdot)}
\newcommand{\uu}{\mathbf{u}}
\newcommand{\vv}{\mathbf{v}}

\renewcommand{\baselinestretch}{1}


\newtheorem{Th}{Теорема}
\newtheorem{Def}{Определение}
\newenvironment{Proof} % имя окружения
{\par\noindent{\textbf{Доказательство.}}} % команды для \begin
{\hfill$\scriptstyle\blacksquare$} % команды для \end
\newtheorem{Assumption}{Предположение}



%\renewcommand{\theequation}{\thesection.\arabic{equation}}
%\renewcommand{\thefigure}{\thesection.\arabic{figure}}
%\renewcommand{\thetable}{\thesection.\arabic{table}}
%\renewcommand{\theequation}{\arabic{section}.\arabic{equation}}
\numberwithin{equation}{section}
%\numberwithin{figure}{section}
%\graphicspath{ {fig/} }

\def\algorithmicrequire{\textbf{\textcolor[rgb]{0.00,0.00,1.00}{Вход:}}}
\def\algorithmicensure{\textbf{\textcolor[rgb]{0.00,0.00,1.00}{Выход:}}}

\usetheme{Warsaw}%{Singapore}%{Warsaw}%{Warsaw}%{Darmstadt}
\usecolortheme{sidebartab}

\setbeamertemplate{footline}[author]
\expandafter\def\expandafter\insertshorttitle\expandafter{\hfill \hfill \hfill \hfill \insertframenumber\,/\,\inserttotalframenumber}

% отключить клавиши навигации
\setbeamertemplate{navigation symbols}{}

\usepackage{wrapfig}
\DeclareCaptionStyle{mystyle}
{format=plain,%
	textformat=period,
	justification=RaggedRight,
	singlelinecheck=true,
}

\DeclareCaptionStyle{singlelineraggedleft}
[justification=RaggedLeft]% right aligned if single line and no `singlelinecheck=false`
{style=mystyle}

\DeclareMathOperator*{\argmax}{argmax}
%\captionsetup{style=singlelineraggedleft}

\makeatletter
\let\@@magyar@captionfix\relax
\makeatother
%----------------------------------------------------------------------------------------------------------


\title[]{Dangers of Bayesian Model Averaging under Covariate Shift}
\author{Pavel Izmailov, Patrick Nicholson, Sanae Lotfi, Andrew Gordon Wilson}

\date{
    2021\,г.
}

\begin{document}

% Creates title page of slide show using above information
\begin{frame}
  \titlepage
\end{frame}

%\begin{frame}
%  \tableofcontents
%\end{frame}

%%%%%%%%%%%%%%%%%%%%%%%%%%%%%%%%%%%%%%%%%%%%%%%%%%%%%%%%%%%%%%%%%%%%%%%%%%%%%%%%%%%%%%%%%%%%%%%%%
%%%%%%%%%%%%%%%%%%%%%%%%%%%%%%%%%%%%%%%%%%%%%%%%%%%%%%%%%%%%%%%%%%%%%%%%%%%%%%%%%%%%%%%%%%%%%%%%%%%%
%\section{Постановка задачи}

%%%%%%%%%%%%%%%%%%%%%%%%%%%%%%%%%%%%%%%%%%%%%%%%%%%%%%%%%%%%%%%%%%%%%%%%%%%%%%%%%%%%%%%%%%%%%%%%%%%%%%
%\subsection{Цель работы}

\begin{frame}{Content}
    \tableofcontents
\end{frame}
\section{Introduction and motivation}

\begin{frame}{The Ubiquitous Problem of Covariate Shift}
    \begin{block}{Real-world machine learning challenge}
        Training and test distributions are rarely identical in practice:
        \[
        p_{\text{train}}(x) \neq p_{\text{test}}(x)
        \]
        but conditional distribution remains: \( p(y|x) \) unchanged.
    \end{block}
    
    \begin{columns}[T]
        \begin{column}{0.48\textwidth}
            \textbf{Applications affected:}
            \begin{itemize}
                \item Medical diagnosis
                \item Autonomous driving
                \item Financial forecasting
                \item Natural language processing
            \end{itemize}
        \end{column}
        \begin{column}{0.48\textwidth}
            \textbf{Types of shift:}
            \begin{itemize}
                \item Test data corruption
                \item Domain shift
                \item Spurious correlations
                \item Temporal drift
            \end{itemize}
        \end{column}
    \end{columns}
    
    \vspace{0.5cm}
    \alert{Key question:} Which methods are robust to such shifts?
\end{frame}


\begin{frame}{Bayesian Methods: Promise vs. Reality}
    \begin{columns}[T]
        \begin{column}{0.48\textwidth}
            \textbf{Theoretical promise:}
            \begin{itemize}
                \item Bayesian Neural Networks (BNNs) provide \textbf{epistemic uncertainty}
                \item Bayesian Model Averaging (BMA):
                \[
                p(y|x) = \int p(y|x,w) p(w|\mathcal{D}) dw
                \]
                \item Should be robust to distribution shift
                \item Captures multiple explanations for OOD data
            \end{itemize}
        \end{column}
        \begin{column}{0.48\textwidth}
            \textbf{Surprising reality:}
            \begin{itemize}
                \item BNNs with high-fidelity HMC inference perform \alert{poorly} under covariate shift
                \item Often \alert{underperform} simple MAP estimation
                \item Particularly vulnerable to noise corruptions
                \item This paper investigates \textbf{why} and proposes solutions
            \end{itemize}
        \end{column}
    \end{columns}
\end{frame}


\begin{frame}{Bayesian neural networks under covariate shift}

    \begin{figure}[!h]
	\centering
	\includegraphics[width=\textwidth]{bmm_1_introduction.jpg}
	\captionsetup{labelformat=empty}
	%\label{fig:graph}
    \end{figure}
    
    
\end{frame}

\section{Theoretical Analysis}

\begin{frame}{Bayesian Framework Review}
    \begin{columns}[T]
        \begin{column}{0.58\textwidth}
            \begin{block}{Bayesian Neural Network}
                \begin{itemize}
                    \item Prior: \( p(w) \)
                    \item Likelihood: \( p(\mathcal{D}|w) \)
                    \item Posterior via Bayes:
                    \[
                    p(w|\mathcal{D}) = \frac{p(\mathcal{D}|w)p(w)}{\int p(\mathcal{D}|w')p(w')dw'}
                    \]
                    \item Prediction via BMA:
                    \[
                    p(y|x) = \int p(y|x,w)p(w|\mathcal{D})dw
                    \]
                \end{itemize}
            \end{block}
        \end{column}
        \begin{column}{0.38\textwidth}
            \begin{block}{MAP Estimation}
                \[
                w_{\text{MAP}} = \arg\max_w p(w|\mathcal{D})
                \]
                Equivalent to:
                \[
                \arg\max_w \left[ \log p(\mathcal{D}|w) + \log p(w) \right]
                \]
                where \(\log p(w)\) acts as regularizer.
            \end{block}
        \end{column}
    \end{columns}
\end{frame}
\begin{frame}
    \begin{theorem}[Lemma 1 - Simplified]
    Consider a BNN with first-layer weight \( w_{ij}^1 \) for input feature \( x^i \). If \( x_k^i = 0 \) for all training examples \( x_k \in \mathcal{D} \), and the prior factorizes:
    \[
    p(W) = p(w_{ij}^1) \cdot p(W \setminus w_{ij}^1)
    \]
    then:
    \[
    p(w_{ij}^1|\mathcal{D}) = p(w_{ij}^1)
    \]
    The posterior equals the prior, and MAP sets \( w_{ij}^1 \) to the prior mode (typically 0).
    \end{theorem}
\end{frame}



\begin{frame}
    \begin{theorem}[Proposition 2 - Key Points]
    Suppose training inputs \( x_1, \ldots, x_n \) lie in an affine subspace:
    \[
    \sum_{j=1}^m x_i^j c_j = c_0 \quad \forall i, \quad \text{with } \sum_{j=0}^m c_j^2 = 1
    \]
    For i.i.d. Gaussian prior \( \mathcal{N}(0,\alpha^2) \):
    \begin{enumerate}
        \item For neuron \( j \) in first layer, posterior of projection \( w_j^c = \sum_{i=1}^m c_i w_{ij}^1 - c_0 b_j^1 \) coincides with prior
        \item MAP sets \( w_j^c = 0 \)
        \item For ReLU activations, BMA predictions are sensitive to test inputs outside the subspace
    \end{enumerate}
    \end{theorem}



\end{frame}
    
\begin{frame}{Another example of covariate shift}
    \begin{figure}[!h]
	\centering
	\includegraphics[width=\textwidth]{bmm_1_robustness.jpg}
	\captionsetup{labelformat=empty}
	%\label{fig:graph}
    \end{figure}

    \begin{figure}[!h]
	\centering
	\includegraphics[width=\textwidth]{bmm_1_robustness_2.jpg}
	\captionsetup{labelformat=empty}
	%\label{fig:graph}
    \end{figure}

\end{frame}

\begin{frame}{Proposed solution}
    \section{Proposed Solution: EmpCov Prior}

The Empirical Covariance (EmpCov) prior for first-layer weights:
\[
p(w^1) = \mathcal{N}\left(0, \alpha\Sigma + \epsilon I\right)
\]
where:
\begin{itemize}
    \item \( \Sigma = \frac{1}{n-1} \sum_{i=1}^n x_i x_i^\top \) is the empirical covariance of inputs
    \item \( \alpha > 0 \) controls prior scale
    \item \( \epsilon > 0 \) ensures positive definiteness
\end{itemize}

\begin{itemize}
    \item \textbf{Eigenvectors of prior} = Principal components of data
    \item \textbf{Prior variance along PC \( p_i \)}: \( \alpha\sigma_i^2 + \epsilon \)
    \item \textbf{For zero-variance direction} (\( \sigma_i^2 = 0 \)): variance = \( \epsilon \) (small)
    \item This prevents BNN from sampling large random weights along low-variance directions
\end{itemize}
    
\end{frame}





\begin{frame}
\begin{table}[H]
\begin{tabular}{lp{0.8\textwidth}}
\toprule
\textbf{Method} & \textbf{Why More Robust to Covariate Shift} \\
\midrule
\textbf{MAP (SGD)} & Regularization (log prior) pushes irrelevant weights to zero \\
\textbf{Deep Ensembles} & Average of multiple MAP solutions, inheriting MAP's robustness \\
\textbf{SWAG} & Gaussian approximation centered around SGD trajectory (near MAP) \\
\textbf{MC Dropout} & Implicit regularization through dropout during training \\
\textbf{Variational Inference} & Often collapses to MAP-like solutions due to approximation \\
\midrule
\textbf{BNN (HMC)} & \textbf{Samples exact posterior, including random weights from prior for low-variance directions} \\
\bottomrule
\end{tabular}
\end{table}

\end{frame}

\begin{frame}
    \begin{figure}[!h]
	\centering
	\includegraphics[width=\textwidth]{bmm_1_final.jpg}
	\captionsetup{labelformat=empty}
	%\label{fig:graph}
    \end{figure}
\end{frame}

    \section{Key Takeaways}

\begin{frame}{Key Takeaways}

\begin{enumerate}
    \item \textbf{BNNs with exact inference (HMC) are vulnerable to covariate shift}, particularly noise corruptions
    \item \textbf{Root cause}: Linear dependencies in training data cause posterior to equal prior in certain directions
    \item \textbf{MAP and ensembles are more robust} due to regularization pushing irrelevant weights to zero
    \item \textbf{EmpCov prior} significantly improves BNN robustness by aligning prior variance with data variance
\end{enumerate}
\end{frame}



%%%%%%%%%%%%%%%%%%%%%%%%%%%%%%%%%%%%%%%%%%%%%%%%%%%%%%%%%%%%%%%%%%%%%%%%%%%





\end{document}
